\section{Derivatives of FPSs}

Our definition of the derivative of a FPS copycats the well-known formula for
the derivative of a power series in analysis:

\begin{definition}
\label{def.fps.deriv}
\lean{AlgebraicCombinatorics.FPS.coeff_derivative_eq}
\leantarget
\leanok
Let $f\in K\left[  \left[  x\right]  \right]  $ be an
FPS. Then, the \emph{derivative} $f^{\prime}$ of $f$ is an FPS defined as
follows: Write $f$ as $f=\sum_{n\in\mathbb{N}}f_{n}x^{n}$ (with $f_{0}%
,f_{1},f_{2},\ldots\in K$), and set%
\[
f^{\prime}:=\sum_{n>0}nf_{n}x^{n-1}.
\]
\end{definition}

\begin{lemma}
\label{lem.fps.deriv.mk}
\lean{AlgebraicCombinatorics.FPS.derivative_eq_mk}
\leanhelper
\leanok
Let $f\in K\left[  \left[  x\right]  \right]  $ be an FPS.
Then $f^{\prime}$ is the power series with coefficient function
$n \mapsto (n+1) \cdot f_{n+1}$.
\end{lemma}

\begin{proof}
\leanok
Immediate from the definition by reindexing $n \mapsto n+1$.
\end{proof}

\begin{lemma}
\label{lem.fps.deriv.eq-derivativeFun}
\lean{AlgebraicCombinatorics.FPS.derivative_eq_derivativeFun}
\leanhelper
\leanok
The derivative operation $f \mapsto f^{\prime}$ agrees with
the standard derivative operation on formal power series.
\end{lemma}

\begin{proof}
\leanok
By definition.
\end{proof}

\begin{lemma}
\label{lem.fps.deriv.coeff-formula}
\lean{AlgebraicCombinatorics.FPS.derivative_coeff_formula}
\leanhelper
\leanok
Let $f\in K\left[  \left[  x\right]  \right]  $ be an FPS.
Then for each $n\in\mathbb{N}$, the $n$-th coefficient of $f^{\prime}$
is $(n+1) \cdot f_{n+1}$.
\end{lemma}

\begin{proof}
\leanok
Follows directly from the definition of the derivative.
\end{proof}

\begin{lemma}
\label{lem.fps.deriv.zero}
\lean{AlgebraicCombinatorics.FPS.derivative_zero}
\leanhelper
\leanok
We have $0^{\prime} = 0$.
\end{lemma}

\begin{proof}
\leanok
All coefficients of $0$ are $0$, so the derivative is $0$.
\end{proof}

\begin{lemma}
\label{lem.fps.deriv.one}
\lean{AlgebraicCombinatorics.FPS.derivative_one}
\leanhelper
\leanok
We have $1^{\prime} = 0$.
\end{lemma}

\begin{proof}
\leanok
$1$ is constant, so $1^{\prime} = 0$.
\end{proof}

\begin{lemma}
\label{lem.fps.deriv.C}
\lean{AlgebraicCombinatorics.FPS.derivative_C}
\leanhelper
\leanok
For any $c\in K$, we have $(c)^{\prime} = 0$
(where $c$ is viewed as a constant FPS).
\end{lemma}

\begin{proof}
\leanok
All coefficients of $c$ beyond the constant term are $0$, so the derivative is $0$.
\end{proof}

\begin{lemma}
\label{lem.fps.deriv.X}
\lean{AlgebraicCombinatorics.FPS.derivative_X}
\leanhelper
\leanok
We have $x^{\prime} = 1$.
\end{lemma}

\begin{proof}
\leanok
The $0$-th coefficient of $x^{\prime}$ is $(0+1) \cdot [x^1]x = 1$, and
all higher coefficients are $0$.
\end{proof}

\begin{lemma}
\label{lem.fps.deriv.X-pow-succ}
\lean{AlgebraicCombinatorics.FPS.derivative_X_pow_succ}
\leanhelper
\leanok
For any $n\in\mathbb{N}$, we have
$\left(x^{n+1}\right)^{\prime} = (n+1) \cdot x^{n}$.
\end{lemma}

\begin{proof}
\leanok
Direct coefficient comparison: $[x^m](x^{n+1})^{\prime} = (m+1) \cdot [x^{m+1}](x^{n+1})$.
This is $(n+1)$ if $m = n$ and $0$ otherwise.
\end{proof}

\begin{lemma}
\label{lem.fps.deriv.X-pow-pos}
\lean{AlgebraicCombinatorics.FPS.derivative_X_pow_of_pos}
\leanhelper
\leanok
For any $n\in\mathbb{N}$ with $n > 0$, we have
$\left(x^{n}\right)^{\prime} = n \cdot x^{n-1}$.
\end{lemma}

\begin{proof}
\leanok
Write $n = m + 1$ and apply Lemma~\ref{lem.fps.deriv.X-pow-succ}.
\end{proof}

\begin{lemma}
\label{lem.fps.deriv.coe-polynomial}
\lean{AlgebraicCombinatorics.FPS.derivative_coe_polynomial}
\leanhelper
\leanok
Let $p$ be a polynomial. Then the derivative of $p$ viewed as a power series
equals the polynomial derivative of $p$ viewed as a power series.
\end{lemma}

\begin{proof}
\leanok
Both derivatives are defined by the same coefficient formula.
\end{proof}

\subsection{Derivative rules}

\begin{theorem}
\label{thm.fps.deriv.rules}
\lean{AlgebraicCombinatorics.FPS.derivative_add, AlgebraicCombinatorics.FPS.derivative_sum, AlgebraicCombinatorics.FPS.derivative_summableFPSSum, AlgebraicCombinatorics.FPS.derivative_smul, AlgebraicCombinatorics.FPS.derivative_mul, AlgebraicCombinatorics.FPS.derivative_div, AlgebraicCombinatorics.FPS.derivative_pow', AlgebraicCombinatorics.FPS.derivative_comp, AlgebraicCombinatorics.FPS.derivative_eq_imp_diff_const}
\leantarget
\leanok

\textbf{(a)} We have $\left(  f+g\right)  ^{\prime
}=f^{\prime}+g^{\prime}$ for any $f,g\in K\left[  \left[  x\right]  \right]
$. \medskip

\textbf{(b)} If $\left(  f_{i}\right)  _{i\in I}$ is a summable family of
FPSs, then the family $\left(  f_{i}^{\prime}\right)  _{i\in I}$ is summable
as well, and we have%
\[
\left(  \sum_{i\in I}f_{i}\right)  ^{\prime}=\sum_{i\in I}f_{i}^{\prime}.
\]


\textbf{(c)} We have $\left(  cf\right)  ^{\prime}=cf^{\prime}$ for any $c\in
K$ and $f\in K\left[  \left[  x\right]  \right]  $. \medskip

\textbf{(d)} We have $\left(  fg\right)  ^{\prime}=f^{\prime}g+fg^{\prime}$
for any $f,g\in K\left[  \left[  x\right]  \right]  $. (This is known as the
\emph{Leibniz rule}.) \medskip

\textbf{(e)} If $f,g\in K\left[  \left[  x\right]  \right]  $ are two FPSs
such that $g$ is invertible, then%
\[
\left(  \dfrac{f}{g}\right)  ^{\prime}=\dfrac{f^{\prime}g-fg^{\prime}}{g^{2}%
}.
\]
(This is known as the \emph{quotient rule}.) \medskip

\textbf{(f)} If $g\in K\left[  \left[  x\right]  \right]  $ is an FPS, then
$\left(  g^{n}\right)  ^{\prime}=ng^{\prime}g^{n-1}$ for any $n\in\mathbb{N}$
(where the expression $ng^{\prime}g^{n-1}$ is to be understood as $0$ if
$n=0$). \medskip

\textbf{(g)} Given two FPSs $f,g\in K\left[  \left[  x\right]  \right]  $, we
have
\[
\left(  f\circ g\right)  ^{\prime}=\left(  f^{\prime}\circ g\right)  \cdot
g^{\prime}%
\]
if $f$ is a polynomial or if $\left[  x^{0}\right]  g=0$. (This is known as
the \emph{chain rule}.) \medskip

\textbf{(h)} If $K$ is a $\mathbb{Q}$-algebra, and if two FPSs $f,g\in
K\left[  \left[  x\right]  \right]  $ satisfy $f^{\prime}=g^{\prime}$, then
$f-g$ is constant.
\end{theorem}

\begin{proof}
This combines parts \textbf{(a)} through \textbf{(h)}, proved individually below.
\end{proof}

\begin{lemma}[Theorem~\ref{thm.fps.deriv.rules} \textbf{(a)}]
\label{lem.fps.deriv.rules.a}
\lean{AlgebraicCombinatorics.FPS.derivative_add}
\leanhelper
\leanok
We have $\left(  f+g\right)  ^{\prime}=f^{\prime}+g^{\prime}$
for any $f,g\in K\left[  \left[  x\right]  \right]$.
\end{lemma}

\begin{proof}
\leanok
This is part of \cite[Exercise 5 \textbf{(b)}]{19s-mt3s} (specifically, it is
Statement 1 in \cite[solution to Exercise 5 \textbf{(b)}]{19s-mt3s}).
\end{proof}

\begin{lemma}
\label{lem.fps.deriv.neg}
\lean{AlgebraicCombinatorics.FPS.derivative_neg}
\leanhelper
\leanok
We have $(-f)^{\prime} = -f^{\prime}$
for any $f\in K\left[  \left[  x\right]  \right]$ (where $K$ is a commutative ring).
\end{lemma}

\begin{proof}
\leanok
The derivative is an additive map, so it preserves negation.
\end{proof}

\begin{lemma}
\label{lem.fps.deriv.sub}
\lean{AlgebraicCombinatorics.FPS.derivative_sub}
\leanhelper
\leanok
We have $(f - g)^{\prime} = f^{\prime} - g^{\prime}$
for any $f,g\in K\left[  \left[  x\right]  \right]$ (where $K$ is a commutative ring).
\end{lemma}

\begin{proof}
\leanok
The derivative is an additive map, so it preserves subtraction.
\end{proof}

\begin{lemma}[Theorem~\ref{thm.fps.deriv.rules} \textbf{(b)}, finite version]
\label{lem.fps.deriv.rules.b.finite}
\lean{AlgebraicCombinatorics.FPS.derivative_sum}
\leanhelper
\leanok
Let $I$ be a finite index set and $(f_i)_{i\in I}$ a family of FPSs.
Then
$\left(\sum_{i\in I} f_i\right)^{\prime} = \sum_{i\in I} f_i^{\prime}$.
\end{lemma}

\begin{proof}
\leanok
Follows from additivity (part \textbf{(a)}) by induction on $|I|$.
\end{proof}

\begin{lemma}
\label{lem.fps.deriv.summable-preserved}
\lean{AlgebraicCombinatorics.FPS.summableFPS_derivative}
\leanhelper
\leanok
If $(f_i)_{i\in I}$ is a summable family of FPSs, then
$(f_i^{\prime})_{i\in I}$ is also summable.
\end{lemma}

\begin{proof}
\leanok
For each coefficient index $n$, the set of $i$ with $[x^n]f_i^{\prime} \neq 0$
is a subset of the set of $i$ with $[x^{n+1}]f_i \neq 0$, which is finite
by summability of $(f_i)$.
\end{proof}

\begin{lemma}[Theorem~\ref{thm.fps.deriv.rules} \textbf{(b)}, infinite version]
\label{lem.fps.deriv.rules.b.infinite}
\lean{AlgebraicCombinatorics.FPS.derivative_summableFPSSum}
\leanhelper
\leanok
If $\left(  f_{i}\right)  _{i\in I}$ is a summable family of
FPSs, then the family $\left(  f_{i}^{\prime}\right)  _{i\in I}$ is summable
as well, and we have%
\[
\left(  \sum_{i\in I}f_{i}\right)  ^{\prime}=\sum_{i\in I}f_{i}^{\prime}.
\]
\end{lemma}

\begin{proof}
This is the natural generalization of Theorem
\ref{thm.fps.deriv.rules} \textbf{(a)} to (potentially) infinite sums. The
proof works coefficient-by-coefficient: for each $n$,
\begin{align*}
[x^n]\left(\sum_{i\in I}f_i\right)^{\prime}
&= \left(\sum_{i\in I}[x^{n+1}]f_i\right) \cdot (n+1) \\
&= \sum_{i\in I} [x^{n+1}]f_i \cdot (n+1)
= \sum_{i\in I} [x^n] f_i^{\prime},
\end{align*}
where the interchange of summation and multiplication by $(n+1)$ is justified
because the sum has finite support.
\end{proof}

\begin{lemma}[Theorem~\ref{thm.fps.deriv.rules} \textbf{(c)}]
\label{lem.fps.deriv.rules.c}
\lean{AlgebraicCombinatorics.FPS.derivative_smul}
\leanhelper
\leanok
We have $\left(  cf\right)  ^{\prime}=cf^{\prime}$ for any $c\in
K$ and $f\in K\left[  \left[  x\right]  \right]  $.
\end{lemma}

\begin{proof}
\leanok
This is part of \cite[Exercise 5 \textbf{(b)}]{19s-mt3s}
(specifically, it is Statement 3 in \cite[solution to Exercise 5 \textbf{(b)}%
]{19s-mt3s}).
\end{proof}

\begin{lemma}
\label{lem.fps.deriv.C-mul}
\lean{AlgebraicCombinatorics.FPS.derivative_C_mul}
\leanhelper
\leanok
We have $(C(c) \cdot f)^{\prime} = C(c) \cdot f^{\prime}$ for any $c\in K$
and $f\in K\left[  \left[  x\right]  \right]$, where $C(c)$ denotes the
constant power series with value $c$.
\end{lemma}

\begin{proof}
\leanok
This follows from the scalar multiplication rule (Lemma~\ref{lem.fps.deriv.rules.c}),
since $C(c) \cdot f = c \cdot f$ in the power series ring.
\end{proof}

\begin{lemma}[Theorem~\ref{thm.fps.deriv.rules} \textbf{(d)}]
\label{lem.fps.deriv.rules.d}
\lean{AlgebraicCombinatorics.FPS.derivative_mul}
\leanhelper
\leanok
We have $\left(  fg\right)  ^{\prime}=f^{\prime}g+fg^{\prime}$
for any $f,g\in K\left[  \left[  x\right]  \right]  $. (This is known as the
\emph{Leibniz rule}.)
\end{lemma}

\begin{proof}
\leanok
This is \cite[Exercise 5 \textbf{(c)}]{19s-mt3s} and
\cite[Proposition 0.2 \textbf{(c)}]{logexp}.
\end{proof}

\begin{lemma}[Theorem~\ref{thm.fps.deriv.rules} \textbf{(e)}]
\label{lem.fps.deriv.rules.e}
\lean{AlgebraicCombinatorics.FPS.derivative_div}
\leanhelper
\leanok
If $f,g\in K\left[  \left[  x\right]  \right]  $ are two FPSs
such that $g$ is invertible, then%
\[
\left(  \dfrac{f}{g}\right)  ^{\prime}=\dfrac{f^{\prime}g-fg^{\prime}}{g^{2}%
}.
\]
(This is known as the \emph{quotient rule}.)
\end{lemma}

\begin{proof}
Let $f,g\in K\left[  \left[  x\right]  \right]  $ be two FPSs
such that $g$ is invertible. Then, Theorem \ref{thm.fps.deriv.rules}
\textbf{(d)} (applied to $\dfrac{f}{g}$ instead of $f$) yields $\left(
\dfrac{f}{g}\cdot g\right)  ^{\prime}=\left(  \dfrac{f}{g}\right)  ^{\prime
}\cdot g+\dfrac{f}{g}\cdot g^{\prime}$. In view of $\dfrac{f}{g}\cdot g=f$,
this rewrites as $f^{\prime}=\left(  \dfrac{f}{g}\right)  ^{\prime}\cdot
g+\dfrac{f}{g}\cdot g^{\prime}$. Solving this for $\left(  \dfrac{f}%
{g}\right)  ^{\prime}$, we find $\left(  \dfrac{f}{g}\right)  ^{\prime}%
=\dfrac{f^{\prime}g-fg^{\prime}}{g^{2}}$.
\end{proof}

\begin{lemma}[Theorem~\ref{thm.fps.deriv.rules} \textbf{(f)}]
\label{lem.fps.deriv.rules.f}
\lean{AlgebraicCombinatorics.FPS.derivative_pow'}
\leanhelper
\leanok
If $g\in K\left[  \left[  x\right]  \right]  $ is an FPS, then
$\left(  g^{n}\right)  ^{\prime}=ng^{\prime}g^{n-1}$ for any $n\in\mathbb{N}$
(where the expression $ng^{\prime}g^{n-1}$ is to be understood as $0$ if
$n=0$).
\end{lemma}

\begin{proof}
\leanok
This follows by induction on $n$, using part \textbf{(d)} (in the
induction step) and $1^{\prime}=0$ (in the induction base).
\end{proof}

\begin{lemma}[Theorem~\ref{thm.fps.deriv.rules} \textbf{(g)}]
\label{lem.fps.deriv.rules.g}
\lean{AlgebraicCombinatorics.FPS.derivative_comp}
\leanhelper
\leanok
Given two FPSs $f,g\in K\left[  \left[  x\right]  \right]  $, we
have
\[
\left(  f\circ g\right)  ^{\prime}=\left(  f^{\prime}\circ g\right)  \cdot
g^{\prime}%
\]
if $f$ is a polynomial or if $\left[  x^{0}\right]  g=0$. (This is known as
the \emph{chain rule}.)
\end{lemma}

\begin{proof}
Let $f,g\in K\left[  \left[  x\right]  \right]  $ be two FPSs
such that $f$ is a polynomial or $\left[  x^{0}\right]  g=0$. Write the FPS
$f$ in the form $f=\sum_{n\in\mathbb{N}}f_{n}x^{n}$ with $f_{0},f_{1}%
,f_{2},\ldots\in K$. Then, $f\left[  g\right]  =\sum_{n\in\mathbb{N}}%
f_{n}g^{n}$. Hence,%
\begin{align}
\left(  f\left[  g\right]  \right)  ^{\prime}  &  =\left(  \sum_{n\in
\mathbb{N}}f_{n}g^{n}\right)  ^{\prime}=\sum_{n\in\mathbb{N}}%
\underbrace{\left(  f_{n}g^{n}\right)  ^{\prime}}_{\substack{=f_{n}\left(
g^{n}\right)  ^{\prime}\\\text{(by part \textbf{(c)})}}}
\ \ \ \ \ \ \ \ \ \ \left(  \text{by part \textbf{(b)}}\right) \nonumber\\
&  =\sum_{n\in\mathbb{N}}f_{n}\left(  g^{n}\right)  ^{\prime}=f_{0}%
\underbrace{\left(  g^{0}\right)  ^{\prime}}_{=1^{\prime}=0}+\sum_{n>0}%
f_{n}\underbrace{\left(  g^{n}\right)  ^{\prime}}_{\substack{=ng^{\prime
}g^{n-1}\\\text{(by part \textbf{(f)})}%
}}=\sum_{n>0}f_{n} n g^{\prime} g^{n-1}. \label{pf.thm.fps.deriv.rules.g.3}%
\end{align}

On the other hand, from $f=\sum_{n\in\mathbb{N}}f_{n}x^{n}$, we obtain
$f^{\prime}=\sum_{n>0}nf_{n}x^{n-1}=\sum_{m\in\mathbb{N}}\left(  m+1\right)
f_{m+1}x^{m}$ (here, we have substituted $m+1$ for $n$ in the sum). Hence,
\[
f^{\prime}\left[  g\right]  =\sum_{m\in\mathbb{N}}\left(  m+1\right)
f_{m+1}g^{m}=\sum_{n>0}nf_{n}g^{n-1}%
\]
(here, we have substituted $n-1$ for $m$ in the sum). Multiplying both sides
of this equality by $g^{\prime}$, we find%
\[
f^{\prime}\left[  g\right]  \cdot g^{\prime}=\left(  \sum_{n>0}nf_{n}%
g^{n-1}\right)  \cdot g^{\prime}=\sum_{n>0}nf_{n}g^{n-1}g^{\prime}.
\]
Comparing this with (\ref{pf.thm.fps.deriv.rules.g.3}), we find $\left(
f\left[  g\right]  \right)  ^{\prime}=f^{\prime}\left[  g\right]  \cdot
g^{\prime}$. In other words, $\left(  f\circ g\right)  ^{\prime}=\left(
f^{\prime}\circ g\right)  \cdot g^{\prime}$ (since $f\circ g$ is a synonym for
$f\left[  g\right]  $).
\end{proof}

\begin{lemma}[Theorem~\ref{thm.fps.deriv.rules} \textbf{(h)}]
\label{lem.fps.deriv.rules.h}
\lean{AlgebraicCombinatorics.FPS.derivative_eq_imp_diff_const}
\leanhelper
\leanok
If $K$ is a $\mathbb{Q}$-algebra, and if two FPSs $f,g\in
K\left[  \left[  x\right]  \right]  $ satisfy $f^{\prime}=g^{\prime}$, then
$f-g$ is constant.
\end{lemma}

\begin{proof}
\leanok
The proof can be found in \cite[Lemma 0.3]{logexp}.
Let $n \geq 1$. From $f^{\prime} = g^{\prime}$, comparing coefficients at index
$n - 1$ yields $n \cdot [x^n]f = n \cdot [x^n]g$. Since $K$ is a $\mathbb{Q}$-algebra
(hence torsion-free), we can divide by $n$ to conclude $[x^n]f = [x^n]g$,
i.e., $[x^n](f - g) = 0$.

Note that the condition that $K$ be a $\mathbb{Q}$-algebra is crucial, since
it allows dividing by positive integers. (For example, if $K$ could be
$\mathbb{Z}/2$, then we would easily get a counterexample, e.g., by taking
$f=x^{2}$ and $g=0$.)
\end{proof}

\begin{lemma}
\label{lem.fps.deriv.eq-of-deriv-eq-const-eq}
\lean{AlgebraicCombinatorics.FPS.eq_of_derivative_eq_of_constantCoeff_eq}
\leanhelper
\leanok
If $K$ is a $\mathbb{Q}$-algebra, and if two FPSs $f,g\in
K\left[  \left[  x\right]  \right]  $ satisfy $f^{\prime}=g^{\prime}$ and
$[x^0]f = [x^0]g$, then $f = g$.
\end{lemma}

\begin{proof}
\leanok
By Lemma~\ref{lem.fps.deriv.rules.h}, $f - g$ is constant. Since
$[x^0](f - g) = [x^0]f - [x^0]g = 0$, we have $f - g = 0$, so $f = g$.
\end{proof}
